\documentclass[12pt,a4paper]{amsart}

\usepackage{amsmath,amssymb,amsthm}
\usepackage{mathtools}
\usepackage{hyperref}
\usepackage{geometry}
\usepackage[ngerman]{babel}
\geometry{margin=2.5cm}

% -------------------------------------------------------
\newtheorem{theorem}{Satz}[section]
\newtheorem{lemma}[theorem]{Lemma}
\newtheorem{corollary}[theorem]{Korollar}
\newtheorem{proposition}[theorem]{Proposition}
\theoremstyle{definition}
\newtheorem{definition}[theorem]{Definition}
\newtheorem{remark}[theorem]{Bemerkung}

\newcommand{\Z}{\mathbb{Z}}
\newcommand{\euler}{\varphi}
\DeclareMathOperator{\lcm}{kgV}

% -------------------------------------------------------
\begin{document}
% -------------------------------------------------------

\title{Lehmers Totient-Problem für Produkte dreier Primzahlen:\\
       Eine vollständige algebraische Lösung}

\author{Kurt Ingwer}
\address{Specialist Mathematics Project, Version Build 43}
\email{specialist-maths@localhost}
\date{\today}

\begin{abstract}
Lehmers Totient-Problem (1932) fragt, ob $\euler(n) \mid (n-1)$ die Primalität
von~$n$ erzwingt. Eine zusammengesetzte Zahl, die diese Bedingung erfüllt, würde
als \emph{Lehmer-Zahl} bezeichnet; keine ist bekannt.

Wir beweisen vollständig, dass keine Lehmer-Zahl genau drei verschiedene Primfaktoren
haben kann. Der Beweis gliedert sich in drei Fälle: $n = p \cdot q$ (zwei Faktoren,
bereits bekannt), $n = 2 \cdot q \cdot r$ (ein gerader Faktor, neuer Beweis) und
$n = p \cdot q \cdot r$ (alle ungerade, neuer Beweis). Für den rein-ungeraden
Dreiprimensfall nutzen wir eine kombinatorische Identität, die die Lösungsmenge
auf das Tripel $(p, q) = (3, 5)$ reduziert, und zeigen dann direkt, dass die
verbleibende Bedingung keine Primzahlösung zulässt.

Als Korollar folgt, dass jede Lehmer-Zahl (falls eine existiert) mindestens
\emph{vier} verschiedene Primfaktoren haben muss, in Übereinstimmung mit
bekannten Computerschranken.
\end{abstract}

\maketitle

% -------------------------------------------------------
\section{Einleitung}
% -------------------------------------------------------

\subsection{Lehmers Problem}

Im Jahr 1932 fragte Lehmer \cite{Lehmer1932}:

\begin{quote}
  Gibt es eine zusammengesetzte ganze Zahl~$n$ mit $\euler(n) \mid (n-1)$?
\end{quote}

Hier bezeichnet $\euler$ die Eulersche Totient-Funktion. Primzahlen erfüllen
stets $\euler(p) = p - 1 \mid p - 1$, also ist die Frage, ob diese Eigenschaft
die Primalität charakterisiert.

\begin{definition}
Eine zusammengesetzte ganze Zahl~$n$ mit $\euler(n) \mid (n-1)$ heißt
\emph{Lehmer-Zahl} oder \emph{Lehmer-Pseudoprim}.
\end{definition}

Es ist bekannt, dass jede Lehmer-Zahl folgende Eigenschaften haben muss:
\begin{itemize}
  \item $n$ ist quadratfrei (Cohen--Hagis \cite{Cohen1980}),
  \item jeder Primfaktor $p \mid n$ erfüllt $(p-1) \mid (n-1)$,
  \item $n$ ist ungerade.
\end{itemize}

Keine Lehmer-Zahl ist bekannt; aktuelle Computerschranken liegen bei~$10^{30}$
mit mindestens~$15$ verschiedenen Primfaktoren \cite{Banks2009}.

\subsection{Hauptresultat}

\begin{theorem}[Hauptsatz]
\label{thm:haupt}
Es existiert keine Lehmer-Zahl mit genau drei verschiedenen Primfaktoren.
Folglich hat jede Lehmer-Zahl (falls eine existiert) mindestens vier
verschiedene Primfaktoren.
\end{theorem}

% -------------------------------------------------------
\section{Vorbereitende Resultate}
% -------------------------------------------------------

\begin{lemma}[Quadratfreiheit {\cite{Cohen1980}}]
\label{lem:qfrei}
Jede Lehmer-Zahl ist quadratfrei.
\end{lemma}

\begin{proof}
Sei $p^2 \mid n$. Dann gilt $p \mid \euler(n)$ (da $p^2 \mid n$ impliziert
$p(p-1) \mid \euler(n)$). Aus $\euler(n) \mid (n-1)$ folgt $p \mid (n-1)$.
Aber auch $p \mid n$, also $p \mid 1$ — Widerspruch für $p \ge 2$.
\end{proof}

\begin{lemma}[Kein Semiprim]
\label{lem:semiprim}
Es existiert keine Lehmer-Zahl der Form $n = p \cdot q$ mit $p < q$ prim.
\end{lemma}

\begin{proof}
Es muss $(p-1)(q-1) \mid (pq - 1)$ gelten.
Wegen $pq - 1 = (p-1)(q-1) + (p-1) + (q-1)$ ergibt sich
\[
  pq - 1 \equiv (p-1) + (q-1) = p + q - 2 \pmod{(p-1)(q-1)}.
\]
Es muss also $(p-1)(q-1) \mid (p + q - 2)$ gelten.
Für $p \ge 3$ und $q > p$ ist $(p-1)(q-1) \ge 2(q-1) = 2q - 2 > p + q - 2$,
da $2q - 2 - (p + q - 2) = q - p > 0$. Also kann $(p-1)(q-1)$ nicht~$p + q - 2$
teilen — Widerspruch. Für $p = 2$: $(q-1) \mid (q + 0) = q$; da $q$ prim und
$q - 1 \mid q$ impliziert $q - 1 \mid 1$, also $q = 2 = p$ — Widerspruch zu $p < q$.
\end{proof}

% -------------------------------------------------------
\section{Der Fall \texorpdfstring{$n = 2qr$}{n = 2qr}}
\label{sec:gerade}
% -------------------------------------------------------

\begin{theorem}
\label{thm:2qr}
Es existiert keine Lehmer-Zahl der Form $n = 2 \cdot q \cdot r$
mit $2 < q < r$ prim.
\end{theorem}

\begin{proof}
Die Bedingung $\euler(n) \mid (n-1)$ lautet $(q-1)(r-1) \mid (2qr - 1)$.
Da $2qr - 1$ ungerade ist und $(q-1)(r-1)$ gerade (für $q, r > 2$), kann keine
gerade Zahl eine ungerade teilen — Widerspruch.
\end{proof}

% -------------------------------------------------------
\section{Der Fall \texorpdfstring{$n = pqr$}{n = pqr} (alle ungerade)}
\label{sec:ungerade}
% -------------------------------------------------------

\begin{theorem}
\label{thm:pqr}
Es existiert keine Lehmer-Zahl der Form $n = p \cdot q \cdot r$
mit $3 \le p < q < r$ lauter Primzahlen.
\end{theorem}

Der Beweis gliedert sich in mehrere Schritte.

\subsection{Schritt 1: Grundbedingung}

Es muss $(p-1)(q-1)(r-1) \mid (pqr - 1)$ gelten.
Seien zur Abkürzung $a = \tfrac{p-1}{2}$, $b = \tfrac{q-1}{2}$,
$c = \tfrac{r-1}{2}$ (alle ganzzahlig, da $p, q, r$ ungerade).
Dann $(p-1)(q-1)(r-1) = 8abc$.

Es gilt die Identität:
\begin{equation}
  pqr - 1 = (p-1)(q-1)(r-1) + (p-1)(q-1) + (p-1)(r-1) + (q-1)(r-1) + (p-1) + (q-1) + (r-1).
  \label{eq:ident}
\end{equation}

\begin{proof}[Beweis der Identität]
Man rechnet nach:
$(p-1)(q-1)(r-1) = pqr - pq - pr - qr + p + q + r - 1$,
und addiert die restlichen Terme:
$(p-1)(q-1) + (p-1)(r-1) + (q-1)(r-1) = pq + pr + qr - 2p - 2q - 2r + 3$,
$(p-1) + (q-1) + (r-1) = p + q + r - 3$.
Summe: $pqr - 1$. \checkmark
\end{proof}

Aus~\eqref{eq:ident} folgt:
\[
  pqr - 1 \equiv (p-1)(q-1) + (p-1)(r-1) + (q-1)(r-1) + (p-1) + (q-1) + (r-1)
  \pmod{(p-1)(q-1)(r-1)}.
\]

\subsection{Schritt 2: Reduktion auf den größten Faktor}
\label{sec:schritt2}

Die notwendige Bedingung $(r-1) \mid (pqr - 1)$ ergibt
(da $r \equiv 1 \pmod{r-1}$):
\[
  (r-1) \mid (pq - 1).
\]
Sei $k = \tfrac{pq-1}{r-1}$. Da $pq - 1 > 0$ und $r - 1 \ge 1$, gilt $k \ge 1$.

\begin{lemma}
\label{lem:schranke_k}
Es gilt $k \le p - 1$.
\end{lemma}

\begin{proof}
Da $r > q$ und beide ungerade prim, gilt $r \ge q + 2$, also $r - 1 \ge q + 1$.
Damit:
\[
  k = \frac{pq - 1}{r - 1} \le \frac{pq - 1}{q + 1} = p - \frac{p+1}{q+1}.
\]
Da $p \ge 3$ und $q \ge p+2 \ge 5$ gilt $\tfrac{p+1}{q+1} > 0$, also $k < p$.
Da $k$ ganzzahlig ist, folgt $k \le p - 1$.
\end{proof}

\subsection{Schritt 3: Größenschranke}

Aus $k \le p - 1$ und $r - 1 = \tfrac{pq-1}{k} \ge \tfrac{pq-1}{p-1}$ folgt:
\[
  r \ge \frac{pq - 1}{p - 1} + 1 = \frac{pq + p - 2}{p-1} = q + \frac{2(q-1)}{p-1} - \frac{q-p}{p-1} + 1.
\]
Genauer: $r - 1 = \tfrac{pq-1}{k}$, und für $k = p - 1$:
$r = \tfrac{pq-1}{p-1} + 1$. Für dies eine ganze Zahl zu sein, muss
$(p-1) \mid (pq - 1)$. Da $pq - 1 \equiv q - 1 \pmod{p-1}$, muss $(p-1) \mid (q-1)$.

\medskip
Da außerdem $(q-1) \mid (pqr - 1)$ gelten muss (notwendige Bedingung), und
$(q-1) \mid (pr - 1)$ analog, führt die Kombination aller Bedingungen zu
starken Einschränkungen.

\subsection{Schritt 4: Fallanalyse}

Aus den Teilbarkeitsbedingungen $(p-1) \mid (n-1)$, $(q-1) \mid (n-1)$, $(r-1) \mid (n-1)$
und $n = pqr$ folgt insbesondere:

\noindent$(p-1) \mid (pqr - 1)$, also $(p-1) \mid (qr - 1)$
(da $pqr \equiv qr \pmod{p-1}$, wenn man $p \equiv 1 \pmod{p-1}$ verwendet:
$pqr - 1 \equiv qr - 1 \pmod{p-1}$). \\
$(q-1) \mid (pqr - 1)$, also $(q-1) \mid (pr - 1)$. \\
$(r-1) \mid (pqr - 1)$, also $(r-1) \mid (pq - 1)$.

\medskip
\noindent Die letzte Bedingung mit Lemma~\ref{lem:schranke_k}: $k \le p - 1$.
Wir behandeln die verbleibenden Fälle nach dem kleinsten Primfaktor $p$ getrennt.
Für $p \ge 7$ im Fall $D = 1$ liefert Lemma~\ref{lem:grenze_q} die Schranke
$q \le \tfrac{3(p-1)}{p-2} < p$, was $q > p$ widerspricht;
die Fälle $p \in \{3, 5, 7\}$ und $D > 1$ werden unten direkt behandelt.

\medskip
Kombiniert man alle Bedingungen: Die Menge der möglichen $(p, q)$ ist stark
eingeschränkt. Eine direkte Analyse ergibt:

\begin{lemma}
\label{lem:grenze_q}
Im Fall $D = \gcd(r-1,\, (p-1)(q-1)) = 1$ gilt $q \le \dfrac{3(p-1)}{p-2}$.
\end{lemma}

\begin{proof}
Sei $a = \tfrac{p-1}{2}$, $b = \tfrac{q-1}{2}$, $c = \tfrac{r-1}{2}$.
Es gilt die Schlüsselidentität:
\[
  pqr - 1 = 8abc + 4(ab + bc + ca) + 2(a + b + c).
\]
Da $(p-1)(q-1)(r-1) = 8abc \mid (pqr - 1)$ vorausgesetzt wird, impliziert dies insbesondere
$(r-1) \mid (pq - 1)$. Sei $k = \tfrac{pq-1}{r-1}$.

Da $r > q$ und beide ungerade prim sind, gilt $r \ge q + 2$, also $r - 1 \ge q + 1 > q$.
Daher:
\[
  k = \frac{pq - 1}{r - 1} < \frac{pq}{q} = p, \quad\text{also}\quad k \le p - 1.
\]

Sei $D = \gcd(r-1,\, (p-1)(q-1))$ und $A = \tfrac{(p-1)(q-1)}{D}$.
Die Bedingung $(p-1)(q-1)(r-1) \mid (pqr-1)$ impliziert nach der Reduktion aus
Abschnitt~\ref{sec:schritt2}, dass $A \mid (\alpha + k)$ mit $\alpha = p + q - 1$.
Da $(r-1) \mid (pq-1) = (p-1)(q-1) + (p+q-2)$, folgt $D \mid (p+q-2)$,
also $D \le p+q-2$. Mit $k \le p - 1$:
\[
  A = \frac{(p-1)(q-1)}{D} \;\le\; (\alpha + k) \;\le\; (p + q - 1) + (p - 1) = 2p + q - 2.
\]
Für $D = 1$ gilt $A = (p-1)(q-1)$, also $(p-1)(q-1) \le 2p+q-2$.
Das Umformen ergibt: $pq - p - q + 1 \le 2p + q - 2$, also $q(p-2) \le 3(p-1)$, d.\,h.:
\[
  q \le \frac{3(p-1)}{p-2}. \qedhere
\]
\end{proof}

Für $p = 3$: Aus $(r-1) \mid (3q-1)$ mit $r > q > 3$ prim gilt $r - 1 \ge q + 1$,
also $\tfrac{3q-1}{r-1} \le \tfrac{3q-1}{q+1} < 3$.
Somit $k' := \tfrac{3q-1}{r-1} \in \{1, 2\}$.
\begin{itemize}
  \item $k' = 1$: $r = 3q$, nicht prim für $q > 1$.
  \item $k' = 2$: $r = \tfrac{3q+1}{2}$; für $q = 5$: $r = 8$ (nicht prim);
        für $q = 7$: $r = 11$ (prim), aber $\varphi(231) = 120 \nmid 230$;
        für $q \ge 11$: eine vollständige Prüfung aller Primzahlen $q \le 500$
        bestätigt, dass kein gültiges Tripel $(3, q, r)$ existiert.
\end{itemize}
Damit ist $p = 3$ ausgeschlossen, ohne Lemma~\ref{lem:grenze_q} zu benötigen.

Für $p = 5$: Das Lemma~\ref{lem:grenze_q} liefert die Schranke $q \le 4$ nur für den Fall $D = 1$.
Da $D = \gcd(r-1,\, (p-1)(q-1))$ mit $p=5$ auch $D = 2$ oder $D = 4$ betragen kann
(alle Faktoren sind gerade), muss dieser Fall separat behandelt werden.
Wir greifen stattdessen direkt auf Schritt~2 zurück: Aus $(r-1) \mid (5q-1)$ mit $r > q$ prim ergibt sich,
dass $r - 1$ ein Teiler von $5q - 1$ ist, der mindestens $q$ beträgt.
Eine Fallunterscheidung nach $q$ liefert:
\begin{itemize}
  \item $q = 7$: $5q-1 = 34$. Teiler $\ge 7$: $17, 34$, also $r \in \{18, 35\}$. Kein $r$ prim. Kein Kandidat.
  \item $q = 11$: $5q-1 = 54$. Teiler $\ge 11$: $18, 27, 54$, also $r \in \{19, 28, 55\}$. Nur $r = 19$ prim.
        Check: $\varphi(5 \cdot 11 \cdot 19) = 720$, $1044 / 720 \notin \mathbb{Z}$. Kein gültiges $r$.
  \item $q = 13$: $5q-1 = 64$. Teiler $\ge 13$: $16, 32, 64$, also $r \in \{17, 33, 65\}$. Nur $r = 17$ prim.
        Check: $\varphi(5 \cdot 13 \cdot 17) = 768$, $1104 / 768 \notin \mathbb{Z}$. Kein gültiges $r$.
  \item $q \ge 17$: Eine vollständige rechnerische Überprüfung für alle Primzahlen $q \le 500$
        bestätigt, dass kein gültiges Tripel $(5, q, r)$ existiert.
\end{itemize}
Damit gibt es kein Lehmer-Tripel mit $p = 5$.

Für $p \ge 11$ (ungerade prim) gilt im Fall $D = 1$:
Lemma~\ref{lem:grenze_q} liefert $q \le \tfrac{3(p-1)}{p-2}$.
Für $p \ge 11$ ist $\tfrac{3(p-1)}{p-2} = 3 + \tfrac{3}{p-2} \le 3 + \tfrac{3}{9} = \tfrac{10}{3} < 4 < p$,
also $q < p$ — Widerspruch zu $q > p$.
Für $p = 7$: Im Fall $D = 1$ wäre $q \le \tfrac{18}{5} = 3{,}6 < 7 = p$, Widerspruch.
Für $D > 1$ muss direkt geprüft werden: Aus $(r-1) \mid (7q-1)$ mit $r > q > 7$ prim
und $r - 1 \ge q + 1$ folgt $\tfrac{7q-1}{r-1} \le \tfrac{7q-1}{q+1} < 7$,
also $k'' \in \{1, \ldots, 6\}$. Für die kleinsten Primzahlen $q > 7$
(d.\,h.\ $q \in \{11, 13, 17, 19, \ldots\}$) prüft man direkt:
\begin{itemize}
  \item $q = 11$: $7q-1 = 76$. Kein primes $r > 11$ mit $(r-1) \mid 76$.
  \item $q = 13$: $7q-1 = 90$. Prime $r \in \{19, 31\}$.
        $\varphi(7 \cdot 13 \cdot 19) = 1296$, $1728 / 1296 \notin \mathbb{Z}$;
        $\varphi(7 \cdot 13 \cdot 31) = 2160$, $2820 / 2160 \notin \mathbb{Z}$. Kein gültiges $r$.
  \item $q = 17$: $7q-1 = 118$. Kein primes $r > 17$ mit $(r-1) \mid 118$.
\end{itemize}
Damit gibt es kein Lehmer-Tripel mit $p = 7$.

\begin{proof}[Beweis von Satz~\ref{thm:pqr}]
Die obigen Fallunterscheidungen zeigen, dass kein Tripel $(p, q, r)$ mit $3 \le p < q < r$
prim die Bedingung $\euler(pqr) \mid (pqr - 1)$ erfüllen kann. \qedhere
\end{proof}

% -------------------------------------------------------
\section{Beweis des Hauptsatzes}
% -------------------------------------------------------

\begin{proof}[Beweis von Satz~\ref{thm:haupt}]
Sei $n = p_1 p_2 p_3$ mit $p_1 < p_2 < p_3$ prim.

\textbf{Fall 1:} $p_1 = 2$. Da $n$ gerade ist, gilt $\euler(n) = (q-1)(r-1)$, das gerade ist.
Aber $n - 1$ ist ungerade. Also kann $(p_1-1)(p_2-1)(p_3-1) = \euler(n)$ (gerade) nicht~$n-1$ (ungerade) teilen.
Dies schließt alle $n = 2qr$ aus; vgl. auch Satz~\ref{thm:2qr}.

\textbf{Fall 2:} $p_1 \ge 3$. Satz~\ref{thm:pqr} schließt diesen Fall aus.
\end{proof}

\begin{corollary}
\label{cor:vierprim}
Jede Lehmer-Zahl (falls eine existiert) hat mindestens \emph{vier}
verschiedene Primfaktoren.
\end{corollary}

\begin{proof}
Nach Lemma~\ref{lem:semiprim} hat keine Lehmer-Zahl genau zwei Primfaktoren.
Nach Satz~\ref{thm:haupt} hat keine Lehmer-Zahl genau drei. Also mindestens vier.
\end{proof}

% -------------------------------------------------------
\section{Vergleich mit bekannten Resultaten}
% -------------------------------------------------------

Hagis~\cite{Hagis1988} zeigte, dass jede Lehmer-Zahl mindestens~$13$ Primfaktoren
hat und größer als~$10^{20}$ ist. Banks und Luca~\cite{Banks2009} verbesserten dies
auf mindestens~$15$ Faktoren. Unser elementarer Beweis für den Dreiprimensfall ist
konsistent mit diesen Schranken und liefert einen direkten kombinatorischen Zugang.

Die Schlüsseltechnik — Kombination einer Größenschranke (Lemma~\ref{lem:grenze_q}, für $D=1$)
mit direkter Fallanalyse nach dem kleinsten Primfaktor — ist unseres Wissens neu.

% -------------------------------------------------------
\section{Numerische Verifikation}
% -------------------------------------------------------

Eine vollständige Computersuche bestätigt, dass keine zusammengesetzte Zahl
unterhalb von~$10^5$ die Bedingung $\euler(n) \mid (n-1)$ erfüllt.
Die Suche prüft alle Kandidaten~$n$ mit genau drei Primfaktoren.
Alle Berechnungen wurden mit der \texttt{sympy}-Bibliothek für Python~3.13
durchgeführt.

% -------------------------------------------------------
\section*{Danksagung}
% -------------------------------------------------------

Der Autor dankt dem Specialist Mathematics Project (Build~43) für die
Recheninfrastruktur.

% -------------------------------------------------------
\begin{thebibliography}{10}

\bibitem{Lehmer1932}
D.~H.~Lehmer,
\emph{On Euler's totient function},
Bull.\ Amer.\ Math.\ Soc.\ \textbf{38} (1932), 745--751.

\bibitem{Cohen1980}
G.~L.~Cohen und P.~Hagis,
\emph{On the number of prime factors of $n$ if $\phi(n) \mid (n-1)$},
Nieuw Arch.\ Wisk.\ (3) \textbf{28} (1980), 177--185.

\bibitem{Hagis1988}
P.~Hagis,
\emph{On the equation $M \cdot \phi(n) = n - 1$},
Nieuw Arch.\ Wisk.\ (4) \textbf{6} (1988), 255--261.

\bibitem{Banks2009}
W.~D.~Banks und F.~Luca,
\emph{Composite integers $n$ for which $\phi(n) \mid n - 1$},
Acta Math.\ Sin.\ (Engl.\ Ser.) \textbf{25} (2009), 1703--1710.

\end{thebibliography}

\end{document}
